\section{Related Work}

\subsection{Self-Repair and Iterative Refinement}

LLM-based self-repair has emerged as a promising paradigm for improving code generation accuracy. \citet{howmanytries2026} systematically evaluated self-repair on HumanEval, finding improvements of 4.9--17.1 percentage points over single-shot generation. They observed that repair success varies by error type, with assertion errors among the more challenging categories. CodeCoR \citep{codecor2025} achieved 77.13\% Pass@1 through multi-agent collaboration, representing current state-of-the-art. However, these works focus on \emph{whether} repair works rather than \emph{why} certain feedback enables better repair.

\citet{debuggingdecay2025} identified ``debugging decay''---substantial capability loss after multiple repair attempts---suggesting that feedback quality in early attempts is crucial. This motivates our focus on single-attempt success and understanding what makes initial feedback effective.

\subsection{Trace-Based Debugging}

DebugRepair \citep{debugrepair2026} demonstrated that runtime traces outperform error messages for code repair, using debugging-style execution to provide richer context. Self-Debug \citep{selfdebug2023} introduced execution trace methodology for self-repair. While these works establish the superiority of traces over messages, they do not decompose \emph{which aspects} of traces drive improvement. Our Actionable Specificity framework attempts this decomposition.

\subsection{Type-Guided Generation}

TyFlow \citep{tyflow2025} showed that type system internalization improves functional correctness by providing structured constraints. This suggests that specific, actionable information (types) helps LLMs generate correct code. Our AS framework generalizes this insight: type information contributes to AS\_loc (type error locations) and AS\_state (type constraints as state).

\subsection{Feedback Informativeness Gap}

Prior work treats feedback signals as categorical (trace vs. message vs. static analysis) rather than decomposing their information content. CoTran \citep{cotran2023} combined compiler and symbolic execution feedback but did not measure which components drive improvement. Meta's Agentic APR \citep{agenticapr2025} achieved 42.3\% solve rate with 11.8 average iterations, demonstrating that iteration quantity alone is insufficient.

\textbf{Our contribution} fills this gap by proposing measurable AS components (localization, state exposure, causal context) and testing their extractability. Our negative result---that regex extraction fails on assertion-dominated benchmarks---reveals a boundary condition that prior work implicitly assumed away.
